\documentclass[12pt]{article}
\usepackage{amsmath,amssymb}

\begin{document}
\section*{{2023 AIME II Problems/Problem 9}}
Source: AoPS Wiki

\subsection*{Solution}
Denote by $O_1$ and $O_2$ the centers of $\omega_1$ and $\omega_2$ , respectively.
Let $XY$ and $AO_1$ intersect at point $C$ .
Let $XY$ and $BO_2$ intersect at point $D$ . Because $AB$ is tangent to circle $\omega_1$ , $O_1 A \perp AB$ .
Because $XY \parallel AB$ , $O_1 A \perp XY$ .
Because $X$ and $P$ are on $\omega_1$ , $O_1A$ is the perpendicular bisector of $XP$ .
Thus, $PC = \frac{PX}{2} = 5$ . Analogously, we can show that $PD = \frac{PY}{2} = 7$ . Thus, $CD = CP + PD = 12$ .
Because $O_1 A \perp CD$ , $O_1 A \perp AB$ , $O_2 B \perp CD$ , $O_2 B \perp AB$ , $ABDC$ is a rectangle. Hence, $AB = CD = 12$ . Let $QP$ and $AB$ meet at point $M$ .
Thus, $M$ is the midpoint of $AB$ .
Thus, $AM = \frac{AB}{2} = 6$ . This is the case because $PQ$ is the radical axis of the two circles, and the powers with respect to each circle must be equal. In $\omega_1$ , for the tangent $MA$ and the secant $MPQ$ , following from the power of a point, we have $MA^2 = MP \cdot MQ$ .
By solving this equation, we get $MP = 4$ . We notice that $AMPC$ is a right trapezoid.
Hence, \begin{align*} AC & = \sqrt{MP^2 - \left( AM - CP \right)^2} \\ & = \sqrt{15} . \end{align*} Therefore, \begin{align*} [XABY] & = \frac{1}{2} \left( AB + XY \right) AC \\ & = \frac{1}{2} \left( 12 + 24 \right) \sqrt{15} \\ & = 18 \sqrt{15}. \end{align*} Therefore, the answer is $18 + 15 = \boxed{\textbf{(033)}}$ . ~Steven Chen (Professor Chen Education Palace, www.professorchenedu.com)

\end{document}
